\documentclass[10pt,a4paper]{amsart}
\usepackage[margin=19mm]{geometry}
\usepackage{amsmath,amssymb,mathtools}
\usepackage[scr=rsfs,cal=euler]{mathalfa}
\usepackage{xfrac}
\usepackage[unicode,hidelinks,bookmarks=false]{hyperref}
\newcommand{\ie}{i.e.\ }
\newcommand{\cf}{cf.\ }
\newcommand{\resp}{respectively\ }
\newcommand{\Subsection}{\subsection}
\providecommand{\nobreakdash}{\nobreak\hbox{-}}
\setlength{\emergencystretch}{2em}
\multlinegap0pt
\usepackage{amssymb}
\usepackage{mathtools}
\newtheorem{lemma}{Lemma}
\newcommand{\m}{\mathbb{M}^n}
\newcommand{\sn}{\mathbb{S}^{n-1}}
\title{Isoperimetric Inequalities and Spectral Consequences in warped product manifolds}
\author{Avas Banerjee}
\date{Source: arXiv:2603.22069v1 (CC BY 4.0)}
\begin{document}
\maketitle

\noindent\emph{Source and licence.} Isoperimetric Inequalities and Spectral Consequences in warped product manifolds. Version arXiv:2603.22069v1 (CC BY 4.0), distributed by arXiv under the Creative Commons Attribution 4.0 International licence, http://creativecommons.org/licenses/by/4.0/, as recorded in the arXiv licence metadata for this exact version and shown on the abstract page https://arxiv.org/abs/2603.22069v1. Full licence notice: \texttt{licenses/arxiv-2603.22069-CC-BY-4.0.txt}.\par\smallskip

\noindent\textit{Editorial scope.} Counted unit: the article's lemma lemma (source0.tex lines 872--877). The article's complete original proof of it is delivered with it (source0.tex lines 879--937, one \texttt{proof} environment of 59 lines). No mathematics is written, repaired or re-derived: the statement and the proof are the article's own lines, in the article's own notation, and no retained line is substituted or re-flowed.  No further source-local context is needed: every cross-reference of every retained line resolves inside the two delivered spans.  Nothing was omitted from any retained span, so the package carries no omission record.  No retained line contains a citation, so the wrapper carries no bibliography and no bibliography entry.  No retained line contains a figure environment, an image-inclusion command or drawing code, so no image is delivered and the package declares no asset dependency.  Editorial formatting only, in addition: an AMS \texttt{amsart} preamble that restates the article's own declarations and macros used by the retained lines, namely the environments \texttt{lemma} on separate counters, together with the standard packages those lines need.  The retained environments print in this document as Lemma 1 in order of appearance, whereas the article prints the same items with the numbering of its own full document.  The complete native source of the article as distributed in the arXiv e-print for this exact version is bundled unchanged as \texttt{sources/196995/source0.tex}.
\begin{lemma}
    Let $\m$ be an n-dimensional manifold with warping function $\psi$. The perimeter of a bounded radial graph set with the associated function $u \in C^1(\sn)$ is given by 
    \begin{equation*}
        \mathrm{Per}_{\m}(\Omega)=\int_{\sn}(\psi(u))^{n-1}\sqrt{1+\frac{|\nabla_{\sn}u|^2}{(\psi(u))^2}}\ \mathrm{d}S.
    \end{equation*}
\end{lemma}

\begin{proof}
    Let $h$ denote the standard metric on $\sn$, with local components
    \begin{equation*}
        h=h_{ij}(\theta)d\theta^id\theta^j,\quad h^{ij}=(h_{ij})^{-1}.
    \end{equation*}
    Since $u$ is $C^1$, the hypersurface $\partial \Omega$ is $C^1$ and $\Omega$ is a set of finite perimeter. Hence
    \begin{equation*}
        \text{Per}_{M}(\Omega)=\mathcal{H}^{n-1}(\partial \Omega).
    \end{equation*}
    We define the parameterization 
    \begin{equation*}
        F: \sn\to \m, \quad F(\theta)=(u(\theta),\theta).
    \end{equation*}
    In local coordinates $(\theta_1,\cdots, \theta^{n-1})$ the tangent vectors to $\partial \Omega$ 
    \begin{equation*}
        X_i:=\partial_i F=\frac{\partial u}{\partial \theta_i}\partial_r+\partial_i,\quad i=1,\cdots, (n-1),
    \end{equation*}
    where $\partial_r$ is the radial coordinate vector and $\partial_i$  are the coordinate vector fields on $\sn$.
Using the structure of the metric $g$,
\begin{equation*}
    g(\partial_r, \partial_r)=1,\quad g(\partial_r, \partial_i)=0,\quad g(\partial_i,\partial_j)=\psi(r)^2h_{ij}.
\end{equation*}
We compute the induced metric $g^{\partial \Omega}$ on $\partial \Omega$: 
\begin{equation*}
    g_{ij}^{\partial \Omega}= g(X_i, X_j)= u_iu_j+\psi(u(\theta))^2h_{ij}(\theta).
\end{equation*}
Let $A=\psi(u(\theta))^2h_{ij}(\theta)$ and $v_{,i}=u_i$ then
\begin{equation*}
    g_{ij}^{\partial \Omega}=A+vv^{T}
\end{equation*}
Since $A$ is positive definite, we apply the matrix determinant identity
\begin{equation*}
    \det(A+vv^T)=\det(A)(1+u^TA^{-1}u).
\end{equation*}
   Now,
   \begin{equation*}
     \det(A)= \psi(u)^{2(n-1)} \det(h_{ij}),\quad A^{-1}=\psi(u)^{-2}h^{ij}.
   \end{equation*}
Therefore,
\begin{equation*}
    \det(g_{ij}^{\partial \Omega})=\psi(u)^{2(n-1)} \det(h_{ij})\left(1+\frac{h^{ij}u_{,i}u_{,j}}{\psi(u)^2}\right).
\end{equation*}
Taking square roots,
\begin{equation*}
    \sqrt{\det(g_{ij}^{\partial \Omega})}= \psi(u)^{n-1}\sqrt{\det(h_{ij})}\sqrt{1+\frac{|\nabla_{\sn}u|^2}{\psi(u)^2}},
\end{equation*}
where
\begin{equation*}
    |\nabla_{\sn}u|^2:= h^{ij}u_{,i}u_{,j}.
\end{equation*}
Since 
\begin{equation*}
    dS(\theta):=\sqrt{\det(h_{ij})}\ \mathrm{d}\theta
\end{equation*}
is the standard surface measure on $\sn$, the $n-1$-dimensional Hausdorff measure on $\partial \Omega$ is
\begin{equation*}
    \text{Per}_{\m}(\Omega)=\mathcal{H}^{n-1}(\partial \Omega)= \int_{\sn}(\psi(u(\theta)))^{n-1}\sqrt{1+\frac{|\nabla_{\sn}u(\theta)|^2}{(\psi(u(\theta)))^2}}\ \mathrm{d}S(\theta).
\end{equation*}
\end{proof}

\end{document}
